\documentclass[10pt,letterpaper]{article}
\usepackage[T1]{fontenc}
\usepackage[utf8]{inputenc}
\usepackage{lmodern}
\usepackage[margin=1in]{geometry}
\usepackage{microtype,graphicx,booktabs,longtable,array}
\usepackage[hidelinks,breaklinks=true]{hyperref}
\newcolumntype{L}[1]{>{\raggedright\arraybackslash}p{#1}}
\newcolumntype{R}[1]{>{\raggedleft\arraybackslash}p{#1}}
\setlength{\parindent}{0pt}
\setlength{\parskip}{5pt}
\setlength{\emergencystretch}{2em}
\renewcommand{\arraystretch}{1.15}
\setlength{\LTpre}{6pt}
\setlength{\LTpost}{6pt}
\hypersetup{pdftitle={@@NAME: SS4A safety analysis evidence},pdfauthor={NJ High Injury Network Generator}}
\begin{document}
\begin{center}
{\Large\bfseries Road Safety Analysis and SS4A Funding Evidence}\par
\vspace{4pt}
{\large @@NAME}\par
@@COUNTY County, New Jersey\par
Observation period: @@START--@@END
\end{center}
\vspace{4pt}
\textbf{Purpose.} This technical report supplies evidence for the safety-analysis
and initial project-development portions of a Safe Streets and Roads for All
(SS4A) submission. It is an input to a locally developed Comprehensive Safety
Action Plan; it is not a completed plan and does not establish funding eligibility.
The program references in this report concern the FY2026 guidance. The applicable
solicitation and local plan must be checked before submission.\cite{nofo,actionplan}

\section{Safety problem and study scope}
The observation period covers @@START--@@END. The loaded, locatable records contain
@@TOTAL crashes. Recorded person totals are @@KILLED killed and @@INJURED injured.
``Unknown'' or ``Not available'' is not zero. Fatal crashes count incidents;
fatalities count people. Injury totals are not counts of seriously injured people.

\subsection{Draft problem statement for local review}
The following paragraph can be adapted for a safety narrative after checking the
source coverage, local geography and reported figures:
\begin{quote}
@@NARRATIVE
\end{quote}

\subsection{Interpretation for SS4A}
The current algorithm selects roadway segments using \emph{all recorded crash
counts}. The FY2026 SS4A definition of a High Injury Network concerns concentrations
of fatal and serious-injury motor-vehicle crashes.\cite{nofo} The network in this
report is therefore described as a \emph{screening network}; it must not be presented
as a verified fatal-and-serious-injury network. There are @@UNKNOWN_SEVERITY records
whose injury detail is unknown in this dataset. Such records cannot be reclassified
as either serious or minor injuries without further source evidence.

The study covers loaded roads assigned to the municipal boundary. It does not
establish municipal ownership, roadway jurisdiction, traffic volumes, population
exposure, or completeness of the jurisdiction's road inventory.

\clearpage
\section{Crash burden and data coverage}
\subsection{Annual crash history}
@@ANNUAL
Annual counts describe the records loaded for each year. A zero loaded count does
not demonstrate a crash-free year; available records do not demonstrate that every
month or archive record is present. Changes between years may reflect coverage or
location processing as well as underlying conditions. No causal or exposure-adjusted
trend is inferred here.

\subsection{Crash severity}
@@SEVERITY
Recorded serious-injury crashes are not the number of seriously injured people.
NJDOT Accidents archives do not consistently provide the serious/minor injury
detail needed for a complete fatal-and-serious-injury analysis. Reconcile the
unknown category with validated source records before using a serious-injury total
in a funding narrative.
\newpage
\subsection{Road-user involvement}
@@USERS
Pedestrian and bicycle involvement categories may overlap. People killed or injured
are subsets of all people, not additional totals to sum. Recorded bicycle ``yes''
counts are not a complete bicycle crash count when involvement is unknown for other
records. Bicycle casualty totals are not inferred from involvement flags.

\subsection{Location quality and completeness}
@@LOCATIONS
Only crashes with usable stored locations are represented. Route-milepost locations
are estimates using the loaded roadway network. Successfully locating a record is
different from assigning it to a road segment: unmatched crashes remain in the
municipal totals but do not contribute to segment rates.

The current analysis does not retain a complete, period-specific accounting of
source records rejected before loading, duplicate removals, records excluded outside
the boundary, or missing source archives. Those quantities must be obtained from
the applicable ingestion reports and original archives. They are not reported as
zero. Preserve those source reports with the application working papers.

\clearpage
\section{Network concentration and corridor evidence}
\subsection{Network concentration}
@@NETWORK
``Eligible loaded roadway miles'' excludes fragments shorter than 0.01 mile, as in
the screening calculation. It is not the mileage of all public roads or of roads
owned by the municipality. The two crash shares above have different denominators:
\emph{all loaded crashes} includes unmatched crashes; \emph{assigned crashes} includes
only records with a road assignment. Each crash contributes once to a selected
segment. Selected miles are based on unique segment identifiers, not a sum of
overlapping corridor buffers.

\begin{figure}[!ht]
\centering
\includegraphics[width=0.8\linewidth]{network.pdf}
\caption{Selected all-crash screening network and loaded roadway context.
Numbered locations identify the ranked segments in Appendix~\ref{segments}.
Unnumbered selected segments may also be shown.}
\end{figure}

\clearpage
\subsection{Corridor evidence}
@@CORRIDORS
``Fatal'' and ``Serious'' count recorded crash incidents. Corridor totals and those
severity counts are saved analysis results; injury-detail-unknown counts come from
the corresponding loaded, assigned records. ``Rate'' means assigned
crashes per segment mile per year, calculated across the corridor. This is not an
exposure-adjusted risk rate. Stored connected same-route groups are used where
available; an ungrouped segment is kept separate. Names and road ownership require
local verification. The table is ordered by crash count to aid review, not by a
benefit-cost score or a finalized investment priority.

For each location being considered, compare the recorded pattern with the original
crash reports, roadway conditions, traffic speeds and volumes, crossing movements,
and pedestrian/bicycle facilities. Document the site visit and identify the road
owner before selecting a treatment. This analysis does not establish which design
change will reduce crashes or by how much.

\section{Community context}
@@EQUITY
The stored overlay uses CDC/ATSDR SVI 2020 national tract information where loaded.
The current join assigns the first intersecting tract ordered by tract identifier;
it is not a length-weighted corridor average. An unknown stored flag remains unknown.
The selected-mile figure above the 75th percentile is geographic context, not an
estimate of the share of crash victims with a particular demographic characteristic.
SVI 2020 is a snapshot, not a contemporaneous measure for every crash year.

Local review should identify schools, transit stops, community destinations,
accessibility barriers and other relevant conditions with dated sources. Those
features are not inferred from crash locations or tract scores in this report.
Record who was consulted, what concerns they raised and how those findings changed
the proposed priorities. No engagement or community support is asserted here.

\clearpage
\section{Material for the SS4A application team}
\subsection{Action Plan components}
Table~\ref{components} relates this output to the seven components in the FY2026
USDOT Action Plan guidance.\cite{actionplan} Local documents, adoption status and the
applicable eligibility worksheet still require review.
\begin{small}
\begin{longtable}{@{}L{100pt}L{144pt}L{182pt}@{}}
\caption{Evidence supplied and local material still needed.}\label{components}\\
\toprule Component & Contribution of this report & Local material still required \\\midrule\endfirsthead
\toprule Component & Contribution of this report & Local material still required \\\midrule\endhead
\bottomrule\endfoot
Leadership commitment and goal setting & Dated safety evidence; no adopted commitment. & Official commitment to eliminate roadway fatalities and serious injuries, with an adopted target date. \\
Planning structure & Reproducible analysis record. & Accountable committee or task force and responsibility for development, implementation and monitoring. \\
Safety analysis & Loaded crash history, screening network, corridor evidence and source limitations. & Validate fatal/serious-injury records, coverage, location patterns and relevant local conditions. \\
Engagement and collaboration & Candidate locations to discuss. & Engagement record, stakeholder participation, interagency coordination and documented effect on priorities. \\
Policy and process changes & No policy audit is performed. & Review existing standards, policies and procedures; identify changes and implementation steps. \\
Strategy and project selections & Initial analytical screening. & Evidence-linked treatments, transparent prioritization, responsible owners and short-, medium- and long-term timeframes. \\
Progress and transparency & A baseline that can be reproduced. & Performance measures, annual accessible public reporting and online publication of the plan. \\
\end{longtable}
\end{small}

\newpage
\subsection{Project development worksheet}
Complete one worksheet for each candidate project selected after local review.
``To be completed locally'' is an unfilled field, not evidence that a requirement
has been met. Connect the project location and proposed scope to a documented safety
problem.\cite{implementation}
\begin{small}
\begin{longtable}{@{}L{140pt}L{290pt}@{}}
\caption{Project information to develop and verify locally.}\\
\toprule Item & Local entry / supporting document \\\midrule\endfirsthead
\toprule Item & Local entry / supporting document \\\midrule\endhead
\bottomrule\endfoot
Project name and limits & To be completed locally: road, endpoints, length, map and referenced segment IDs. \\
Sponsor and road owner & To be completed locally: responsible entity, jurisdiction and any required agreements. \\
Safety problem & Adapt the factual narrative and cite the relevant annual, corridor and crash-source evidence. Confirm severity detail. \\
Proposed scope & To be completed locally: engineering-reviewed treatment, supporting design criteria and relationship to the observed pattern. \\
Project selection & To be completed locally: prioritization criteria, alternatives, community input and adopted-plan reference. \\
Design and feasibility & To be completed locally: design stage, field review, permits, environmental review, right-of-way and utility needs. \\
Cost estimate & To be completed locally: itemized quantities, unit costs, contingencies, estimate date and preparation basis. No costs are generated here. \\
Funding and match & To be completed locally: eligible project cost, requested federal share, non-federal sources and commitments, checked against the applicable notice. \\
Delivery schedule & To be completed locally: design, approvals, procurement and construction milestones with accountable parties. \\
Operations and maintenance & To be completed locally: owner, responsibilities and recurring resource needs. \\
Expected results & To be completed locally: evidence-supported benefits and evaluation method. No crash-reduction or benefit-cost estimate is generated here. \\
Monitoring & To be completed locally: baseline period, comparable follow-up data, responsible party and public reporting schedule. \\
\end{longtable}
\end{small}

The FY2026 NOFO limits the federal share to 80 percent of eligible costs.\cite{nofo}
This is a dated reference, not an eligibility or match calculation. The FY2026
application period closed on May 26, 2026; confirm the applicable solicitation
before preparing a request. Do not assume a future round uses the same rules.

\subsection{Application-specific fatality measures}
USDOT's FY2026 application instructions use 2019--2023 motor-vehicle roadway
fatalities for the average annual fatality-rate calculation.\cite{fatalityrate}
This report's user-selected period and loaded local records are not automatically
that measure. Verify the specified geography, fatality source and population
denominator separately. The app has not calculated the SS4A application fatality
rate, a benefit-cost ratio or a complete serious-injury person total.

\clearpage
\section{Methods and reproducibility}
Loaded crashes with usable locations are assigned to the nearest road segment
within @@SNAP meters. Segment rates use actual assigned crash counts divided by
segment miles and the number of selected years. A one-sided Poisson screening test
uses a leave-one-out road-class reference rate, then a municipality fallback.
Selection requires at least three crashes and a p-value below @@THRESHOLD; no
reference exposure means no selection. Fragments shorter than 0.01 mile are excluded
from screening and reference exposure. Same-route endpoints within one meter form
connected corridor groups. Severity weighting is a separate descriptive measure,
not the Poisson event count or the ranking used in this report.

Traffic-volume exposure, spatial dependence, over-dispersion and multiple testing
are not modeled. This is exploratory screening and requires transportation-safety
review. Crash records come from loaded NJDOT county/year Accidents archives;
road geometry comes from the loaded NJDOT measured network and municipal boundaries
from NJ Office of GIS. Source completeness and network vintage must be checked in
the ingestion records.

\textbf{Analysis ID:} @@ANALYSIS_ID. \textbf{Method:} @@METHOD.\par
\textbf{Run created:} @@RUN_DATE. \textbf{Report generated:} @@GENERATED.
@@VERSIONS

\clearpage
\appendix
\section{Ranked segment detail}\label{segments}
The map numbers refer to this table, which is ordered by crash rate rather than the
corridor table's total crash count. Rates are recorded crashes per mile per year.
Only the top ten segments are listed; retain the full CSV/GeoJSON export as a
working attachment when developing a project.
@@SEGMENTS

\begin{thebibliography}{9}
\raggedright
\bibitem{actionplan} U.S. Department of Transportation. \emph{Comprehensive Safety Action Plans.}
Updated March 27, 2026. \url{https://www.transportation.gov/grants/ss4a/comprehensive-safety-action-plans}.
\bibitem{nofo} U.S. Department of Transportation. \emph{Safe Streets and Roads for All: FY2026 Notice of Funding Opportunity.}
Issued March 27, 2026. \url{https://www.transportation.gov/sites/dot.gov/files/2026-03/SS4A-FY26-NOFO.pdf}.
\bibitem{implementation} U.S. Department of Transportation. \emph{Implementation Grants.}
Updated April 7, 2026. \url{https://www.transportation.gov/grants/ss4a/implementation-grants}.
\bibitem{fatalityrate} U.S. Department of Transportation. \emph{Calculating Average Annual Fatality Rate for SS4A.}
Updated March 27, 2026. \url{https://www.transportation.gov/grants/ss4a/calculating-average-annual-fatality-rate-ss4a}.
\end{thebibliography}
\end{document}
